
Contrastive self-supervised learning (SSL) is widely assumed to amplify spurious correlations. Its objective encourages representations where features stable across random crops and color jitter become strongly instance-discriminative. Salient spurious features such as background texture satisfy both criteria --- they are stable across typical augmentations and vary strongly across instances --- so the augmentation-invariance argument predicts that contrastive SSL backbones should encode spurious attributes more strongly than supervised counterparts. This work provides direct empirical evidence that the opposite is true on the Waterbirds benchmark.

In controlled experiments on frozen ResNet-50 representations, supervised ERM encodes spurious background features significantly more strongly than contrastive MoCo-v3 (spurious/task probe ratio: ERM = 1.052, MoCo-v3 = 1.027; difference = 0.025, Cohen's d = 5.68, p\_bonf < 0.0001, Bonferroni-corrected over six pairs). The paradigm ranking on Waterbirds is ERM $\approx$ DINO > BarlowTwins > MoCo-v3 --- the reverse of what augmentation-invariance theory predicts. On CelebA, this ranking reverses, establishing a dataset-dependent interaction between spurious attribute type and pretraining objective. An ablation in which SimCLR is retrained from scratch with background-replacement augmentation shows a 23.1\% reduction in spurious/task ratio relative to standard augmentation (p = 1.26 $\times 10^{-6})$, confirming that augmentation design causally modulates spurious feature encoding.

\textbf{The gap this work fills.} Despite substantial prior work on spurious correlations in deep learning [Sagawa et al., 2020; Kirichenko et al., 2022; Geirhos et al., 2020], systematic cross-paradigm comparison of spurious feature encoding in frozen representations is absent. DFR \citep{Kirichenko2022} demonstrates that ERM features are sufficient for robustness after group-balanced retraining but does not compare to SSL. WILDS \citep{Koh2021WILDS} and DomainBed \citep{Gulrajani2021DomainBed} evaluate fine-tuned models, conflating pretraining representation quality with fine-tuning dynamics. To the best of our knowledge, we provide the first controlled four-paradigm comparison of spurious/task probe accuracy ratio on frozen ResNet-50 features across two group-annotated benchmarks, combined with a causal augmentation ablation.

\textbf{Key insight.} The dominant driver of spurious feature encoding is supervised label correlation, not augmentation-based instance-discriminability. ERM's cross-entropy objective is trained on data where spurious attributes reliably predict class labels (Waterbirds: 95\% training correlation), causing the model to jointly encode both task and spurious features. MoCo-v3, being label-agnostic, encodes spurious features only insofar as they are instance-discriminative under the augmentation set --- a weaker pressure that measurably reduces spurious encoding. DINO's self-distillation, which implicitly generates class-level semantic targets via momentum teacher, produces spurious encoding indistinguishable from ERM (d = 0.60, p\_bonf = 1.0), consistent with this interpretation.

\textbf{Contributions.} We make the following contributions:

\begin{enumerate}
\item \textbf{Waterbirds paradigm comparison:} A controlled four-paradigm $\times$ five-seed comparison on frozen ResNet-50 demonstrates a highly significant paradigm effect on spurious/task probe accuracy ratio (ANOVA F = 35.99, p = 2.42 $\times 10^{-7})$. ERM encodes spurious background features most strongly and MoCo-v3 least, opposite to augmentation-invariance theory.
\end{enumerate}

\begin{enumerate}
\item \textbf{Cross-dataset replication:} CelebA replication (ANOVA F = 5.51, p = 0.009) confirms the paradigm effect but reveals a ranking reversal (MoCo-v3 lowest on Waterbirds, highest on CelebA), establishing that spurious attribute type interacts with pretraining objective.
\end{enumerate}

\begin{enumerate}
\item \textbf{Mechanistic interpretation:} Supervised label correlation is identified as the primary driver of spurious feature encoding, supported by the ERM $\approx$ DINO similarity (d = 0.60, p\_bonf = 1.0) and the large ERM > MoCo-v3 effect (d = 5.68).
\end{enumerate}

\begin{enumerate}
\item \textbf{Augmentation ablation:} Background-replacement augmentation in contrastive training produces a 23.1\% reduction in spurious/task ratio relative to standard augmentation (paired t = 46.66, p = 1.26 $\times 10^{-6}$, d = 32.4), confirming augmentation design as a functional causal lever.
\end{enumerate}

